already proved, we conclude that for any @@M0@@,
\sgaStableTarget{sga6.e13.eq.2-9-2}@@M1@@
for every @@M2@@.

By 2.5, @@M3@@ is a @@M4@@-sheaf. Put @@M5@@ for @@M6@@ and @@M7@@. By 1.11, @@M8@@ and @@M9@@ is @@M10@@-regular. Thus, by 1.4 (ii), @@M11@@ for every @@M12@@. Therefore, if we define polynomials by @@M13@@ and @@M14@@, then we have (ii) by 2.9.2.

Finally, since @@M15@@ is @@M16@@-regular, @@M17@@ for @@M18@@ by 1.4 (i)\,; hence @@M19@@. But @@M20@@ by 1.8 (i). Thus, if we define @@M21@@, we have (iii).

\sgaStableTarget{sga6.e13.rem.2-10}\textbf{Remark 2.10.} For fixed @@M22@@, consider an equality of polynomials in @@M23@@:
@@M24@@
Then one easily finds the formula
@@M25@@
by induction on @@M26@@ so that @@M27@@ becomes @@M28@@, and by taking @@M29@@.

Similarly, one finds that there exists a polynomial with rational coefficients @@M30@@ which is increasing in @@M31@@, such that in the following equality:
@@M32@@
